% ==============================================================================
% MACROS, PACKAGES COMPLÉMENTAIRES ET COMMANDES PERSONNALISÉES (ÉDITION COMPACTE)
% ==============================================================================

\usepackage{amsmath,amssymb,amsfonts,amsthm}
\usepackage{mathtools}
\usepackage{bm}
\usepackage{booktabs}
\usepackage{tabularx}
\usepackage{multirow}
\usepackage{makecell}
\usepackage{graphicx}
\usepackage{subcaption}
\usepackage{float}
\usepackage{xcolor}
\usepackage{tikz}
\usetikzlibrary{shapes,arrows.meta,positioning,calc,fit,backgrounds,decorations.pathreplacing}
\usepackage{pgfplots}
\pgfplotsset{compat=1.18}
\usepackage{listings}
\usepackage{tcolorbox}
\tcbuselibrary{skins,breakable}
\usepackage{hyperref}
\usepackage{csquotes}
\usepackage{titlesec}

% Palette de couleurs institutionnelle
\definecolor{EnsamBlue}{RGB}{21, 60, 85}        % Bleu Officiel ENSAM
\definecolor{primaryNavy}{RGB}{15, 23, 42}      % Slate 900
\definecolor{secondaryBlue}{RGB}{30, 58, 138}   % Blue 900
\definecolor{accentIndigo}{RGB}{79, 70, 229}    % Indigo 600
\definecolor{crimsonRisk}{RGB}{185, 28, 28}     % Red 700
\definecolor{amberWarning}{RGB}{217, 119, 6}    % Amber 600
\definecolor{emeraldSuccess}{RGB}{5, 150, 105}  % Emerald 600
\definecolor{lightSlate}{RGB}{248, 250, 252}    % Slate 50
\definecolor{borderSlate}{RGB}{203, 213, 225}   % Slate 300
\definecolor{codeBackground}{RGB}{245, 247, 250}% Code BG
\definecolor{codeComment}{RGB}{100, 116, 139}   % Slate 500
\definecolor{codeString}{RGB}{13, 148, 136}     % Teal 600
\definecolor{codeKeyword}{RGB}{67, 56, 202}     % Indigo 700

% Configuration compacte des titres académiques (Standard Mémoire d'Ingénieur)
\titleformat{\chapter}[hang]
  {\normalfont\Large\bfseries\color{primaryNavy}}
  {\thechapter.\ }
  {0pt}
  {}
\titlespacing*{\chapter}{0pt}{-10pt}{12pt}
\titlespacing*{\section}{0pt}{10pt}{4pt}
\titlespacing*{\subsection}{0pt}{8pt}{3pt}

% Configuration des hyperliens
\hypersetup{
    colorlinks=true,
    linkcolor=secondaryBlue,
    citecolor=accentIndigo,
    urlcolor=secondaryBlue,
    filecolor=crimsonRisk,
    bookmarksnumbered=true,
    pdfstartview={FitH}
}

% Opérateurs mathématiques personnalisés
\DeclareMathOperator*{\argmin}{arg\,min}
\DeclareMathOperator*{\argmax}{arg\,max}
\newcommand{\inDeg}{\operatorname{deg}^-}
\newcommand{\outDeg}{\operatorname{deg}^+}
\newcommand{\cycleID}{\texttt{cycle\_id}}
\newcommand{\alertID}{\texttt{alert\_id}}
\newcommand{\txnID}{\texttt{transaction\_id}}

% Boîtes d'avertissement et de remarques académiques compactes
\newtcolorbox{academicnote}[1][Remarque Académique]{
    colback=lightSlate,
    colframe=secondaryBlue,
    coltitle=white,
    fonttitle=\bfseries\footnotesize,
    title=#1,
    arc=1.5mm,
    boxrule=0.75pt,
    left=2.5mm,
    right=2.5mm,
    top=1.5mm,
    bottom=1.5mm
}

\newtcolorbox{regulatoryalert}[1][Exigence Réglementaire (GAFI / UTRF)]{
    colback=lightSlate,
    colframe=crimsonRisk,
    coltitle=white,
    fonttitle=\bfseries\footnotesize,
    title=#1,
    arc=1.5mm,
    boxrule=0.75pt,
    left=2.5mm,
    right=2.5mm,
    top=1.5mm,
    bottom=1.5mm
}

% Configuration des listings de code source compacte
\lstset{
    backgroundcolor=\color{codeBackground},
    basicstyle=\ttfamily\scriptsize,
    breaklines=true,
    breakatwhitespace=true,
    captionpos=b,
    commentstyle=\color{codeComment}\itshape,
    keywordstyle=\color{codeKeyword}\bfseries,
    stringstyle=\color{codeString},
    numbers=left,
    numberstyle=\tiny\color{codeComment},
    numbersep=5pt,
    frame=single,
    framerule=0.4pt,
    rulecolor=\color{borderSlate},
    tabsize=2,
    showstringspaces=false,
    keepspaces=true
}

\lstdefinelanguage{JSON}{
    keywords={true,false,null},
    morestring=[b]",
    morecomment=[l]{//}
}

\lstdefinelanguage{YAML}{
    keywords={true,false,null,yes,no},
    morestring=[b]",
    morecomment=[l]{\#}
}

% Commande unifiée d'insertion de figure pleine largeur
\newcommand{\singlefigure}[4][0.92\linewidth]{%
    \begin{figure}[htbp]
        \centering
        \includegraphics[width=#1]{#2}%
        \caption{#3}%
        \label{#4}%
    \end{figure}
}

% Commande d'insertion de 2 captures côte-à-côte (hauteur 4.4cm)
% Usage : \screenshotpaired[hauteur]{img1}{caption1}{img2}{caption2}{Main caption}{Main label}
\newcommand{\screenshotpaired}[7][4.4cm]{%
    \begin{figure}[htbp]
        \centering
        \begin{subfigure}[b]{0.485\textwidth}
            \centering
            \includegraphics[width=\linewidth,height=#1,keepaspectratio]{#2}
            \caption{#3}
        \end{subfigure}\hfill
        \begin{subfigure}[b]{0.485\textwidth}
            \centering
            \includegraphics[width=\linewidth,height=#1,keepaspectratio]{#4}
            \caption{#5}
        \end{subfigure}
        \vspace{1mm}
        \caption{#6}
        \label{#7}
    \end{figure}
}
